\section{Graph dynamics: exact identities and limits}
\label{sec:dynamics}

\subsection{Typed message passing}

Assume each of $d$ relation types supplies one incoming neighbor per vertex.
For relation $r$, let $\pi_r$ be a permutation and set
$(P_rx)_i=x_{\pi_r(i)}$. Hence
\begin{equation}
 P_r\1=\1,\qquad P_r^\top\1=\1,\qquad P_r^\top P_r=I.
 \label{eq:permutation}
\end{equation}
Self edges and duplicate neighbors are separately prohibited in the pilot;
those exclusions do not follow merely from \eqref{eq:permutation}.

For the binary alphabet let $s_i=2X_i-1$. With $d=8$, the ordered local
input and two-logit rule are
\begin{align}
 u_i&=(s_i,(P_1s)_i,\ldots,(P_8s)_i)^\top\in\{-1,1\}^{9},\\
 z_i&=W_2\ReLU(W_1u_i+b_1)+b_2,\qquad
 X_i^{t+1}=\argmax_{a\in\{0,1\}}z_{i,a}.
 \label{eq:binaryrule}
\end{align}
Use $W_1\in\R^{16\times9}$ and $W_2\in\R^{2\times16}$.
The $9\cdot16+16+2\cdot16+2=194$ rule parameters are sampled and frozen.
Relation labels are retained; this is not an unordered average of neighbors.

For $q$ categorical states, replace $s_i$ by its $q$-dimensional one-hot
encoding and gather a vector in $\R^{(d+1)q}$. A spatial perception map
$\R^{9q}\to\R^4$, then $\R^4\to\R^{16}\to\R^q$, is one possible
categorical rule family. It is different from the binary family in
\eqref{eq:binaryrule}, even if both use the same graph.

\begin{proposition}[Grid equivalence]
\label{prop:convolution}
On a periodic square grid with the eight Moore-neighbor relation types,
the first affine map in \eqref{eq:binaryrule} is exactly a periodic
$3\times3$ cross-correlation with 16 output channels, after matching
the central/neighborhood order. The remaining maps are pointwise.
\end{proposition}
\begin{proof}
For hidden channel $h$, the gathered expression is
\[
 (W_1u_i+b_1)_h
 =b_{1,h}+\sum_{\delta\in\{-1,0,1\}^2}
                   k_{h,\delta}s_{i+\delta},
\]
where $k_{h,\delta}$ is the corresponding column of $W_1$. Periodic
indexing makes this precisely the cross-correlation expression.
ReLU and the output affine map act at each vertex independently.
Matching columns and offsets therefore gives identical logits and argmax
states, including a common tie convention.
\end{proof}

Re-expressing a fixed grid calculation as message passing therefore
introduces no new mathematical treatment. Rewiring changes the inputs to
the shared rule; it is not, by itself, a faster numerical implementation.

\begin{proposition}[Equivariance under consistent relabeling]
For a vertex permutation $\Pi$, replace $x$ by $\Pi x$ and each $P_r$ by
$\Pi P_r\Pi^\top$. The shared rule then satisfies
\[
 F_{\omega,\Pi G}(\Pi x)=\Pi F_{\omega,G}(x).
\]
\end{proposition}
\begin{proof}
The new gathered source is
$(\Pi P_r\Pi^\top)(\Pi x)=\Pi P_rx$. Every row receives the same local
function, so row permutation commutes with the update.
\end{proof}
This is not invariance of a raster-serialized transformer with fixed
position embeddings: relabeling also changes its input sequence.

\subsection{What endpoint swaps preserve}

Swapping $\pi_r(a)$ and $\pi_r(b)$ leaves $\pi_r$ bijective. Thus it
preserves one incoming and one outgoing edge of type $r$ at every vertex.
Degree alone does not preserve connectivity, distances, eigenvalues,
motifs, rule outputs, or the distribution of serialized token sequences.

An explicit counterexample uses 64 vertices split into two sets of 32.
Within each set, use the eight cyclic offsets
$\{-4,-3,-2,-1,1,2,3,4\}$ modulo 32. Every relation is a permutation;
there are no self or duplicate neighbors and all degrees equal eight.
Yet the graph is disconnected. Degree constraints alone therefore cannot
guarantee better information mixing than the connected periodic grid.
This does not assert that the pilot's particular swap sequence produces
that counterexample.

\subsection{An exactly solvable diffusion surrogate}

To study a limited meaning of mixing, define the linear operator
\begin{equation}
 Q_G=\frac{I+\sum_{r=1}^{8}P_r}{9}.
 \label{eq:diffusion}
\end{equation}
It is doubly stochastic, so the uniform mean is preserved.
\emph{This is a linear surrogate, not the random ReLU--argmax NCA.}
On an $s\times s$ periodic Moore grid, Fourier vectors
$v_{ab}(r,c)=\exp(2\pi\mathrm i(ar+bc)/s)$ are eigenvectors with
\begin{equation}
 \lambda_{ab}=
 \frac{(1+2\cos(2\pi a/s))(1+2\cos(2\pi b/s))}{9}.
 \label{eq:fourier}
\end{equation}
This follows by summing the nine separable shifts.
For $s=8$, the largest nonconstant eigenvalue magnitude is
\begin{equation}
 \rho=\frac{1+\sqrt2}{3}=0.8047378541,\qquad
 1-\rho=\frac{2-\sqrt2}{3}=0.1952621459.
\end{equation}
Indeed, the largest nonconstant product in \eqref{eq:fourier} is
$3(1+\sqrt2)$; all negative products have magnitude at most three.
Since $Q_{\rm grid}$ is symmetric, if $\bar x=n^{-1}\1^\top x$,
\begin{equation}
 \norm{Q_{\rm grid}^t(x-\bar x\1)}_2
 \leq \rho^t\norm{x-\bar x\1}_2.
\end{equation}
Reducing this norm bound by a factor $\delta$ requires
\begin{equation}
 t\geq\left\lceil\frac{\log\delta}{\log\rho}\right\rceil
 =\begin{cases}11,&\delta=0.1,\\22,&\delta=0.01.\end{cases}
\end{equation}
These are exact bound calculations, not simulated relaxation times.

For directed rewiring, a general $Q_G$ need not be symmetric.
The appropriate one-step Euclidean bound on mean-zero signals is
$\norm{Q_G-J}_2$, where $J=\1\1^\top/n$, rather than an unqualified
claim about its second eigenvalue. The disconnected example has a
nonconstant componentwise-constant vector with eigenvalue one.
Nothing here proves that rewiring increases the relevant gap, or that a
larger diffusion gap improves language pre-training.

\subsection{How many local inputs can change?}

Let
\begin{equation}
 h=\sum_{i=1}^n\sum_{r=1}^{d}
       \mathbf1\{\pi_r^G(i)\ne\pi_r^0(i)\}
\end{equation}
count changed \emph{typed source slots}. If
$S_Gx=(x,P_1^Gx,\ldots,P_d^Gx)$ is the stacked gather operator, then
\begin{equation}
 \norm{S_G-S_0}_F^2
 =\sum_{r=1}^d\norm{P_r^G-P_r^0}_F^2=2h.
 \label{eq:gatherfrobenius}
\end{equation}
Each changed row replaces one unit entry by a different unit entry.
For independent binary $X_i$ with probability $1/2$,
$\E[(X_j-X_k)^2]=1/2$ for distinct $j,k$, giving
\begin{equation}
 \E\norm{(S_G-S_0)X}_2^2=\frac h2.
 \label{eq:binaryperturb}
\end{equation}
The preserved pilot reports changed-slot fraction $113/512=0.220703125$,
so $h=113$ and \eqref{eq:binaryperturb} equals $56.5$.
This identity concerns the \emph{first update's common independent initial
state}; later states become correlated. A nominal attempted rewiring rate
of $0.25$ is not the same quantity.

At a common state supplied to both topologies, for a node with $c_i$
changed slots, its signed local vectors
obey $\norm{u_i^G-u_i^0}_2\leq2\sqrt{c_i}$. Let
$m(u)=z_1(u)-z_0(u)$ and
\[
 L_m=\norm{W_{2,1:}-W_{2,0:}}_2\norm{W_1}_2.
\]
ReLU is 1-Lipschitz, so
\begin{equation}
 \abs{m(u_i^G)-m(u_i^0)}\leq2L_m\sqrt{c_i}.
\end{equation}
If $\abs{m(u_i^0)}>2L_m\sqrt{c_i}$, the binary prediction cannot change.
Outside that margin condition it may change. Argmax is discontinuous at
zero margin; a global Lipschitz bound for the discrete NCA must not be
inferred from a bound for its logits.

